%; whizzy paragraph -pdf xpdf -latex ./whizzypdfptex.sh
%; whizzy-paragraph "^\\\\begin{frame}"
% latex beamer presentation.
% platex, latex-beamer でコンパイルすることを想定。 

%     Tokyo Debian Meeting resources
%     Copyright (C) 2009 Junichi Uekawa
%     Copyright (C) 2009 Nobuhiro Iwamatsu

%     This program is free software; you can redistribute it and/or modify
%     it under the terms of the GNU General Public License as published by
%     the Free Software Foundation; either version 2 of the License, or
%     (at your option) any later version.

%     This program is distributed in the hope that it will be useful,
%     but WITHOUT ANY WARRANTY; without even the implied warreanty of
%     MERCHANTABILITY or FITNESS FOR A PARTICULAR PURPOSE.  See the
%     GNU General Public License for more details.

%     You should have received a copy of the GNU General Public License
%     along with this program; if not, write to the Free Software
%     Foundation, Inc., 51 Franklin St, Fifth Floor, Boston, MA  02110-1301 USA

\documentclass[cjk,dvipdfmx,12pt]{beamer}
\usetheme{Tokyo}
\usepackage{monthlypresentation}

%  preview (shell-command (concat "evince " (replace-regexp-in-string "tex$" "pdf"(buffer-file-name)) "&")) 
%  presentation (shell-command (concat "xpdf -fullscreen " (replace-regexp-in-string "tex$" "pdf"(buffer-file-name)) "&"))
%  presentation (shell-command (concat "evince " (replace-regexp-in-string "tex$" "pdf"(buffer-file-name)) "&"))

%http://www.naney.org/diki/dk/hyperref.html
%日本語EUC系環境の時
\AtBeginDvi{\special{pdf:tounicode EUC-UCS2}}
%シフトJIS系環境の時
%\AtBeginDvi{\special{pdf:tounicode 90ms-RKSJ-UCS2}}

\title{Debian JP 会長就任の挨拶}
%\subtitle{第75回 2011年4月度}
\author{まえだこうへい mkouhei@debian.or.jp\\IRC nick: mkouhei}
\date{2011年4月16日}
\logo{\includegraphics[width=8cm]{image200607/openlogo-light.eps}}

\begin{document}

\frame{\titlepage{}}

\section{}

\begin{frame}
 \frametitle{アンタ誰？}
  \begin{itemize}
   \item ネコ(こまめ)一匹、ヨメ一人
   \item Debianを使い始めたのはSargeリリースの直前くらい
   \item OpenBlockS, Armadilloなどで利用
   \item 主に東京エリアDebian勉強会の運営など
   \item KVMの記事(@IT, KVM徹底入門), CouchDB JP
  \end{itemize}
\end{frame}

\begin{frame}{さて問題です。}

{\Huge なにが問題でしょう？}

\end{frame}

\begin{frame}{Debin JP Projectの現状}
 \begin{enumerate}
  \item 若者(20代以下)が少ない
  \item 新たな加入者が少ない
  \item (会長・監事選挙が)マンションの管理組合のようだ
 \end{enumerate}
\end{frame}

\begin{frame}{若者が少ない}
 \begin{itemize}
  \item 身近な人にきっかけ作りを(職場、家族、近所)
  \item 学生との交流を増やしたい
  \item そして若い世代が中心になるようにしたい
 \end{itemize}
 \begin{center}
  {\Huge 鉄は熱いうちに打て}
 \end{center}
\end{frame}

\begin{frame}{新たな加入者が少ない}
 \begin{itemize}
  \item 他の勉強会、コミュニティとの交流
  \item 関心を持ってそうな人の勧誘
  \item 業務で使っている、使う予定の技術者、企業との交流も
 \end{itemize}
 \begin{center}
  {\Huge 情けは人の為ならず}
 \end{center}
\end{frame}

\begin{frame}{マンションの管理組合のようだ}
 \begin{itemize}
  \item 誰もいなかったので推薦をお受けすることに
  \item 「自分は向いてない」「DDじゃないから」は止めよう
  \item (来期以降は)未経験者が自ら立候補するようにしたい
 \end{itemize}
 \begin{center}
  {\Huge 自重・自粛無用}
 \end{center}
\end{frame}

\begin{frame}{既存活動は今までどおり}
 \begin{itemize}
  \item  「Debianは頼られているので、いままでどおりやっていきましょう」by (c)2010 荒木さん
  \item 「Debian開発者の会」なのだから、なんらかの開発（ソフトウェア、パッ
	ケージメンテナンス、翻訳）をやるべし by (c)2009 岩松さん
 \end{itemize}
 \begin{center}
  {\Large と、自分のハードルをあげてみる}
 \end{center}
\end{frame}

\begin{frame}{まとめ}
\begin{itemize}
 \item Debianに興味をもってもらい、ユーザを増やし、そこから将来の開発者を増やしましょう
 \item 既にDebianを利用していても自分たちだけで悶々としている人を仲間に入れましょう
 \item これまでの活動は今までどおり続けていきましょう
\end{itemize}
\end{frame}
\end{document}
